\documentclass[11pt,a4paper]{article}
\usepackage[utf8]{inputenc}
\usepackage[T1]{fontenc}
\usepackage[french]{babel}
\usepackage{geometry}
\usepackage{hyperref}
\usepackage{xcolor}
\usepackage{enumitem}
\usepackage{fancyhdr}
\usepackage{tikz}
\usepackage{tcolorbox}
\usepackage{fontawesome5}
\usetikzlibrary{shapes.geometric,arrows.meta,positioning,fit,backgrounds,decorations.pathreplacing}

\geometry{margin=2.5cm}
\hypersetup{colorlinks=true, linkcolor=blue, urlcolor=blue}

\tcbuselibrary{skins,breakable}
\newtcolorbox{infobox}[1][]{
  colback=blue!5, colframe=blue!40, title=#1, fonttitle=\bfseries\small,
  breakable, left=4pt, right=4pt}
\newtcolorbox{warnbox}[1][]{
  colback=orange!5, colframe=orange!60, title=#1, fonttitle=\bfseries\small,
  breakable, left=4pt, right=4pt}
\newtcolorbox{tipbox}[1][]{
  colback=green!5, colframe=green!50!black, title=#1, fonttitle=\bfseries\small,
  breakable, left=4pt, right=4pt}

\pagestyle{fancy}
\fancyhf{}
\rhead{Actuarial Notebook Agent}
\lhead{Guide Utilisateur}
\rfoot{Page \thepage}

\title{%
  \textbf{Actuarial Notebook Agent}\\[0.5em]
  \large Guide Utilisateur\\[0.3em]
  \normalsize Version 2.0 — Mars 2026
}
\author{}
\date{}

\begin{document}
\maketitle
\tableofcontents
\newpage

% ─────────────────────────────────────────────────────────────────────────────
\section{Introduction}
% ─────────────────────────────────────────────────────────────────────────────

L'\textbf{Actuarial Notebook Agent} vous permet de construire une table de mortalité
d'expérience à partir d'un fichier CSV de portefeuille assuré, sans écrire une seule
ligne de code. L'application se compose de \textbf{3 onglets} :

\medskip
\begin{center}
\begin{tikzpicture}[
  tab/.style={rectangle, draw, rounded corners=4pt, minimum width=3.8cm,
              minimum height=1.2cm, fill=blue!12, font=\small\bfseries},
  arrow/.style={-Stealth, thick, blue!60}
]
  \node[tab] (t1) {📚 Bibliothèque\\{\normalfont\scriptsize Sélectionner les notebooks}};
  \node[tab, right=1.5cm of t1] (t2) {📋 Ordre d'exécution\\{\normalfont\scriptsize Réordonner les étapes}};
  \node[tab, right=1.5cm of t2] (t3) {🔬 Analyse et rapport\\{\normalfont\scriptsize Charger CSV + résultats}};
  \draw[arrow] (t1) -- (t2);
  \draw[arrow] (t2) -- (t3);
\end{tikzpicture}
\end{center}

% ─────────────────────────────────────────────────────────────────────────────
\section{Démarrage rapide}
% ─────────────────────────────────────────────────────────────────────────────

\begin{enumerate}[leftmargin=1.5cm, label=\textbf{Étape \arabic*.}]
  \item \textbf{Lancer l'application} depuis un terminal :
\begin{verbatim}
streamlit run app.py
\end{verbatim}
        L'interface s'ouvre dans votre navigateur à l'adresse
        \texttt{http://localhost:8501}.

  \item \textbf{Onglet 3 — Analyse} : glissez votre fichier CSV dans la zone de dépôt.

  \item L'agent lance automatiquement l'analyse. Chaque étape s'affiche progressivement
        avec son rapport.

  \item À la fin, un \textbf{résumé exécutif} s'affiche et le fichier CSV de la table
        d'expérience est exporté dans le répertoire de travail.
\end{enumerate}

\begin{infobox}[Pré-requis : fichier CSV]
Le fichier doit contenir les colonnes :
\texttt{id}, \texttt{date\_naissance}, \texttt{date\_entree}, \texttt{date\_sortie},
\texttt{cause\_sortie} (\texttt{'deces'} ou \texttt{'autre'}), \texttt{sexe} (\texttt{'H'} ou \texttt{'F'}).
\end{infobox}

% ─────────────────────────────────────────────────────────────────────────────
\section{Onglet 1 — Bibliothèque de notebooks}
% ─────────────────────────────────────────────────────────────────────────────

\subsection{Vue de l'écran}

\begin{center}
\begin{tikzpicture}[font=\small]
  % Fenêtre application
  \draw[rounded corners=6pt, fill=gray!8, draw=gray!50] (0,0) rectangle (14,10);
  \node[anchor=north west, font=\bfseries\normalsize] at (0.3,9.8) {Actuarial Notebook Agent};

  % Tabs
  \foreach \x/\lbl in {1/📚 Bibliothèque, 5/📋 Ordre, 9/🔬 Analyse} {
    \draw[fill=blue!15, draw=blue!40, rounded corners=3pt] (\x,8.8) rectangle (\x+3.5,9.5);
    \node[font=\scriptsize] at (\x+1.75,9.15) {\lbl};
  }
  % Tab active
  \draw[fill=white, draw=blue!60, rounded corners=3pt, line width=1pt] (0.8,8.8) rectangle (4.5,9.5);
  \node[font=\scriptsize\bfseries] at (2.65,9.15) {📚 Bibliothèque};

  % Left panel — checkboxes
  \draw[fill=white, draw=gray!40, rounded corners=3pt] (0.5,0.5) rectangle (5.5,8.5);
  \node[anchor=north west, font=\bfseries\scriptsize] at (0.7,8.4) {Sélectionner les notebooks :};
  \foreach \y/\nb in {7.5/01\_chargement\_donnees.ipynb, 6.5/02\_controle\_qualite.ipynb,
                       5.5/03\_expositions\_taux\_bruts.ipynb, 4.5/04\_lissage\_whittaker\_henderson.ipynb,
                       3.5/05\_visualisation.ipynb, 2.5/06\_smr.ipynb, 1.5/07\_export.ipynb} {
    \draw[fill=blue!20, draw=blue!50] (0.7,\y-0.15) rectangle (1.0,\y+0.15);
    \node[anchor=west, font=\tiny] at (1.1,\y) {\nb};
  }
  \draw[fill=green!20, draw=green!50, rounded corners=2pt] (0.7,0.7) rectangle (2.8,1.2);
  \node[font=\tiny] at (1.75,0.95) {✅ Tout sélectionner};
  \draw[fill=red!10, draw=red!30, rounded corners=2pt] (3.0,0.7) rectangle (5.3,1.2);
  \node[font=\tiny] at (4.15,0.95) {☐ Tout déselectionner};

  % Right panel — preview
  \draw[fill=white, draw=gray!40, rounded corners=3pt] (5.7,0.5) rectangle (13.5,8.5);
  \node[anchor=north west, font=\bfseries\scriptsize] at (5.9,8.4) {Aperçu du contenu :};
  \foreach \y/\lbl in {7.6/📓 01\_chargement\_donnees.ipynb ▶, 6.8/📓 02\_controle\_qualite.ipynb ▶,
                        6.0/📓 03\_expositions\_taux\_bruts.ipynb ▶} {
    \draw[fill=blue!8, draw=blue!30, rounded corners=2pt] (5.9,\y-0.25) rectangle (13.3,\y+0.25);
    \node[anchor=west, font=\tiny] at (6.0,\y) {\lbl};
  }

  % Test button
  \draw[fill=orange!20, draw=orange!60, rounded corners=3pt] (0.7,1.3) rectangle (13.3,1.8);
  \node[font=\scriptsize\bfseries] at (7,1.55) {▶ Exécuter les notebooks sélectionnés (données synthétiques)};
\end{tikzpicture}
\end{center}

\subsection{Actions disponibles}

\begin{description}
  \item[Cocher/décocher un notebook] Inclure ou exclure ce notebook de l'analyse.
        Par défaut, tous les 7 notebooks sont cochés.
  \item[Aperçu] Cliquer sur \og \texttt{📓 nom.ipynb ▶} \fg{} pour voir le contenu
        (code et documentation) du notebook.
  \item[Exécuter le blueprint] Lance l'analyse sur 50 000 contrats synthétiques
        (données calibrées sur la table TH 00-02) pour valider la méthodologie.
\end{description}

\begin{tipbox}[Conseil]
Utilisez le bouton \og Exécuter le blueprint \fg{} pour vérifier que la chaîne fonctionne
avant de charger vos vraies données. Le SMR devrait être proche de 1,0 sur les données synthétiques.
\end{tipbox}

% ─────────────────────────────────────────────────────────────────────────────
\section{Onglet 2 — Ordre d'exécution}
% ─────────────────────────────────────────────────────────────────────────────

\subsection{Vue de l'écran}

\begin{center}
\begin{tikzpicture}[font=\small]
  \draw[rounded corners=6pt, fill=gray!8, draw=gray!50] (0,0) rectangle (14,8);

  % Tabs
  \draw[fill=white, draw=blue!60, rounded corners=3pt, line width=1pt] (4.8,7.3) rectangle (8.5,8.0);
  \node[font=\scriptsize\bfseries] at (6.65,7.65) {📋 Ordre d'exécution};
  \foreach \x/\lbl in {0.5/📚 Bibliothèque, 9/🔬 Analyse} {
    \draw[fill=blue!15, draw=blue!40, rounded corners=3pt] (\x,7.3) rectangle (\x+3.5,8.0);
    \node[font=\scriptsize] at (\x+1.75,7.65) {\lbl};
  }

  % Content
  \node[anchor=north west, font=\bfseries\small] at (0.5,7.0) {Ordre d'exécution des notebooks};

  \foreach \y/\nb/\desc in {
    6.0/01\_chargement\_donnees.ipynb/Étape 1 — Imports et chargement des données,
    5.0/02\_controle\_qualite.ipynb/Étape 2 — Contrôle qualité des données,
    4.0/03\_expositions\_taux\_bruts.ipynb/Étape 3 — Calcul des expositions,
    3.0/04\_lissage\_whittaker\_henderson.ipynb/Étape 4 — Lissage WH,
    2.0/05\_visualisation.ipynb/Étape 5 — Visualisation} {
    % Up button
    \draw[fill=gray!20, draw=gray!50, rounded corners=2pt] (0.5,\y-0.3) rectangle (1.1,\y+0.3);
    \node[font=\scriptsize] at (0.8,\y) {↑};
    % Down button
    \draw[fill=gray!20, draw=gray!50, rounded corners=2pt] (1.2,\y-0.3) rectangle (1.8,\y+0.3);
    \node[font=\scriptsize] at (1.5,\y) {↓};
    % Info
    \node[anchor=west, font=\tiny\bfseries] at (2.0,\y+0.1) {\nb};
    \node[anchor=west, font=\tiny, text=gray] at (2.0,\y-0.15) {\desc};
  }

  % Reset button
  \draw[fill=yellow!20, draw=yellow!60, rounded corners=3pt] (0.5,0.3) rectangle (3.5,0.9);
  \node[font=\scriptsize] at (2.0,0.6) {↺ Réinitialiser l'ordre};

  % Summary
  \draw[fill=blue!10, draw=blue!30, rounded corners=3pt] (0.5,1.1) rectangle (13.5,1.7);
  \node[font=\tiny] at (7.0,1.4) {7 notebooks actifs : 01 → 02 → 03 → 04 → 05 → 06 → 07};
\end{tikzpicture}
\end{center}

\subsection{Comment réordonner}

\begin{itemize}
  \item \textbf{Bouton ↑} : remonte le notebook d'une position.
  \item \textbf{Bouton ↓} : descend le notebook d'une position.
  \item \textbf{↺ Réinitialiser} : remet l'ordre alphabétique d'origine.
\end{itemize}

\begin{warnbox}[Attention à l'ordre]
Les notebooks partagent des variables (ex.\,\texttt{data}, \texttt{table}).
L'ordre par défaut (01 → 07) doit être conservé pour une analyse standard.
Ne réordonner que si vous savez ce que vous faites (ex.\,insérer une étape de correction).
\end{warnbox}

% ─────────────────────────────────────────────────────────────────────────────
\section{Onglet 3 — Analyse et rapport}
% ─────────────────────────────────────────────────────────────────────────────

\subsection{Vue de l'écran — chargement}

\begin{center}
\begin{tikzpicture}[font=\small]
  \draw[rounded corners=6pt, fill=gray!8, draw=gray!50] (0,0) rectangle (14,6);

  \draw[fill=white, draw=blue!60, rounded corners=3pt, line width=1pt] (9.2,5.3) rectangle (13.0,6.0);
  \node[font=\scriptsize\bfseries] at (11.1,5.65) {🔬 Analyse et rapport};
  \foreach \x/\lbl in {0.5/📚 Bibliothèque, 4.8/📋 Ordre} {
    \draw[fill=blue!15, draw=blue!40, rounded corners=3pt] (\x,5.3) rectangle (\x+3.5,6.0);
    \node[font=\scriptsize] at (\x+1.75,5.65) {\lbl};
  }

  \node[anchor=north west, font=\bfseries\small] at (0.5,5.0) {Analyse sur données réelles};

  % Upload zone
  \draw[dashed, draw=blue!50, rounded corners=4pt, fill=blue!5] (0.5,2.5) rectangle (13.5,4.5);
  \node[font=\large] at (7,3.8) {📁};
  \node[font=\small\bfseries] at (7,3.3) {Charger un fichier de portefeuille (CSV)};
  \node[font=\scriptsize, text=gray] at (7,2.9) {Glissez-déposez votre fichier ici ou cliquez pour parcourir};
  \node[font=\scriptsize, text=gray] at (7,2.65) {Formats acceptés : .csv};

  % Info message
  \draw[fill=blue!10, draw=blue!30, rounded corners=3pt] (0.5,0.3) rectangle (13.5,2.2);
  \node[font=\small\bfseries] at (7,1.6) {7 notebooks sélectionnés.};
  \node[font=\small] at (7,1.2) {Chargez un fichier CSV pour démarrer l'analyse automatique.};
\end{tikzpicture}
\end{center}

\subsection{Vue de l'écran — analyse en cours}

\begin{center}
\begin{tikzpicture}[font=\small]
  \draw[rounded corners=6pt, fill=gray!8, draw=gray!50] (0,0) rectangle (14,9);

  % Spinner
  \draw[fill=blue!10, draw=blue!40, rounded corners=3pt] (0.5,8.0) rectangle (13.5,8.7);
  \node[font=\small] at (7,8.35) {⏳ Étape 3 : Calcul des expositions et taux bruts de mortalité};

  % Step 1 results
  \node[anchor=north west, font=\bfseries\small, color=blue!80] at (0.5,7.7) {Étape 1 — Imports et chargement des données};
  \draw[fill=white, draw=gray!30, rounded corners=3pt] (0.5,5.5) rectangle (13.5,7.4);
  \node[anchor=north west, font=\small] at (0.7,7.3) {Les données ont été chargées avec succès :};
  \node[anchor=west, font=\small] at (1.0,6.9) {• \textbf{50 000} contrats analysés};
  \node[anchor=west, font=\small] at (1.0,6.5) {• \textbf{2 200} décès observés (4,4 \%)};
  \node[anchor=west, font=\small] at (1.0,6.1) {• Période 2010–2023, sexe \textbf{H}, paramètres par défaut.};
  \draw[gray!40] (0.5,5.5) -- (13.5,5.5);

  % Step 2 results
  \node[anchor=north west, font=\bfseries\small, color=blue!80] at (0.5,5.3) {Étape 2 — Contrôle qualité};
  \draw[fill=white, draw=gray!30, rounded corners=3pt] (0.5,3.5) rectangle (13.5,5.0);
  \node[anchor=north west, font=\small] at (0.7,4.9) {Aucune anomalie détectée. \textbf{50 000} lignes conservées.};
  \node[anchor=west, font=\small] at (1.0,4.4) {Âge moyen à l'entrée : \textbf{52,3} ans — Durée moy. observation : \textbf{6,1} ans.};
  \draw[gray!40] (0.5,3.5) -- (13.5,3.5);

  % Figure placeholder
  \draw[dashed, fill=gray!5, draw=gray!40, rounded corners=3pt] (0.5,0.3) rectangle (13.5,3.3);
  \node[font=\small\itshape, text=gray] at (7,1.8) {[graphique taux bruts vs lissés s'affichera ici]};
\end{tikzpicture}
\end{center}

\subsection{Vue de l'écran — rapport final}

À la fin de l'analyse, un résumé exécutif s'affiche :

\begin{center}
\begin{tikzpicture}[font=\small]
  \draw[rounded corners=6pt, fill=gray!8, draw=gray!50] (0,0) rectangle (14,5);

  \node[anchor=north west, font=\bfseries\large] at (0.5,4.8) {Synthèse de l'analyse};
  \draw[fill=white, draw=blue!20, rounded corners=3pt] (0.5,0.3) rectangle (13.5,4.5);
  \node[anchor=north west, font=\small] at (0.7,4.3) {\textbf{Portefeuille :} 50 000 contrats hommes, période 2010–2023.};
  \node[anchor=north west, font=\small] at (0.7,3.8) {\textbf{Plage d'âges :} 20–90 ans (53 âges avec données suffisantes).};
  \node[anchor=north west, font=\small] at (0.7,3.3) {\textbf{SMR global :} 0,9987 — mortalité conforme à la table TH 00-02.};
  \node[anchor=north west, font=\small] at (0.7,2.8) {\textbf{Lissage :} Whittaker-Henderson λ=100, courbe monotone.};
  \node[anchor=north west, font=\small] at (0.7,2.3) {\textbf{Fichier exporté :} \texttt{table\_mortalite\_synthetique\_H.csv}};
\end{tikzpicture}
\end{center}

% ─────────────────────────────────────────────────────────────────────────────
\section{Questions fréquentes}
% ─────────────────────────────────────────────────────────────────────────────

\subsection{L'analyse démarre-t-elle automatiquement ?}
Oui. Dès qu'un fichier CSV est chargé, l'agent lance automatiquement les 7 étapes
avec les paramètres par défaut (sexe=H, date\_fin=2023-12-31, λ=100).

\subsection{Puis-je changer le sexe ou la date de fin d'observation ?}
Dans la version actuelle, les paramètres sont définis dans le notebook 01.
L'agent les adapte automatiquement si vous l'indiquez dans le nom du fichier
ou si vous modifiez le notebook 01 avant l'analyse.

\subsection{Le rapport montre-t-il du code Python ?}
Non. Seuls les résultats formatés (texte, tableaux, graphiques) sont affichés.
Le code Python reste dans les notebooks et n'est jamais visible dans le rapport.

\subsection{Que faire si l'agent détecte une erreur ?}
L'agent diagnostique l'erreur, corrige le code et relance l'étape automatiquement.
Si l'erreur persiste après plusieurs tentatives, un message d'erreur rouge s'affiche
avec le traceback.

\subsection{Le fichier CSV de résultats est-il exporté automatiquement ?}
Oui, dans le répertoire de travail, sous le nom
\texttt{table\_mortalite\_<nom\_fichier>\_<sexe>.csv}.

\subsection{Puis-je ajouter mes propres notebooks ?}
Oui. Déposez votre fichier \texttt{.ipynb} dans le dossier \texttt{notebooks/},
rechargez l'application : il apparaît automatiquement dans l'onglet 1.

\begin{infobox}[Rappel sur les variables partagées]
Votre notebook peut utiliser toutes les variables créées par les notebooks précédents
(\texttt{data}, \texttt{table}, \texttt{SEXE}, \texttt{SMR}, etc.).
\end{infobox}

% ─────────────────────────────────────────────────────────────────────────────
\section{Interprétation des résultats}
% ─────────────────────────────────────────────────────────────────────────────

\subsection{Le SMR}

\begin{center}
\begin{tabular}{cll}
\toprule
\textbf{Valeur SMR} & \textbf{Interprétation} & \textbf{Action recommandée} \\
\midrule
$< 0{,}80$ & Très forte sous-mortalité & Vérifier la qualité des données \\
$0{,}80$–$0{,}90$ & Sélection favorable marquée & Normal pour portefeuille épargne \\
$0{,}90$–$1{,}10$ & Conforme à la référence & Table directement utilisable \\
$1{,}10$–$1{,}25$ & Sur-mortalité modérée & Analyser par tranche d'âge \\
$> 1{,}25$ & Sur-mortalité forte & Vérifier composition du portefeuille \\
\bottomrule
\end{tabular}
\end{center}

\subsection{Qualité du lissage}

\begin{itemize}
  \item \textbf{Âges non-monotones $= 0$} : courbe correcte, utilisable directement.
  \item \textbf{Âges non-monotones $> 0$} : augmentez λ (ex.\,λ=200 ou λ=500).
  \item \textbf{Nombreux âges avec $E_x < 10$} : diminuez λ pour moins lisser
        (les données ont peu de crédibilité sur ces âges).
\end{itemize}

\end{document}
